\begin{figure}[H]
\centering
\begin{tikzpicture}[scale=1.12,transform shape,
  >=Latex,
  box/.style={draw,rounded corners,minimum width=2.4cm,minimum height=8mm,
              align=center},
  note/.style={font=\footnotesize,align=center},
  every node/.style={font=\small}
]
  \node[box] (nullity) at (0,1.8)
    {$\dim\Cond(C,C^\perp)$\\is $0$ or $1$};
  \node[box] (torus) at (-3.4,0)
    {$\F_q^2\rtimes T$\\dimension $0$};
  \node[box] (sl) at (3.4,0)
    {$\F_q^2\rtimes\mathrm{SL}_2(q)$\\dimension $1$};
  \draw[->] (nullity) -- node[above left,note] {no diagonal\\duality witness} (torus);
  \draw[->] (nullity) -- node[above right,note] {unique projective\\witness $[s]$} (sl);

  \node[box] (nt) at (-3.4,-2.0)
    {$\F_q^2\rtimes N(T)$\\input blocks across encoder views};
  \draw[->,dashed] (torus) -- node[right,note]
    {realized odd\\coordinate permutation} (nt);

  \node[note] at (3.4,-1.7)
    {the linear $\mathrm{SL}_2(q)$ factor has a Weil lift;\\
     the affine translation factor has a Heisenberg obstruction};
\end{tikzpicture}
\caption{The dimension of the code-to-dual conductor selects the exact
fixed-coordinate projective transversal group.  Coordinate permutations form
a separate extension problem: in the non-GRS examples, an odd
permutation enlarges the possible input blocks from $T$ to $N(T)$ across
encoder views, where $N(T)$ is the normalizer of $T$ in
$\mathrm{SL}_2(q)$; the permutation need not fix the chosen encoder input.}
\label{fig:symmetry-groups}
\end{figure}
